% ---------------------------------------------------------------- APÊNDICE: ANÁLISE COMPLETA DAS PROVAS (para o professor e para quem quiser os detalhes)
\clearpage
\titulo{Análise completa}{Análise completa das provas de 2021 a 2025}\label{pg:analise}
Esta parte mostra os números por trás do material. Serve ao professor e a quem quiser saber de onde vieram as escolhas.

\subsecao{Os 20 assuntos, do que mais cai ao que menos cai}
\begin{center}
\begin{tikzpicture}
\begin{axis}[
  xbar, width=0.70\linewidth, height=12.2cm, bar width=9pt,
  xmin=0, xmax=12.5, ymin=0.3, ymax=20.7,
  axis x line=none, axis y line*=left, axis line style={cinza!60}, tick style={draw=none},
  ytick={1,...,20}, y dir=reverse,
  yticklabels={\textbf{Circuitos elétricos e Lei de Ohm},\textbf{Cinemática (movimento)},\textbf{Leis de Newton e forças},\textbf{Acústica (som)},
    Calorimetria e mudanças de fase,Propagação de calor,Ondas: propriedades e fenômenos,Ondas eletromagnéticas e radiações,Eletromagnetismo,
    {Energia, trabalho e potência},Potência e energia elétrica,Termodinâmica e gases,Óptica geométrica,Física moderna,
    Hidrostática,Gravitação,Eletrostática,Termometria e dilatação,Impulso e quantidade de movimento,Estática e equilíbrio},
  yticklabel style={font=\small,color=tinta},
  nodes near coords, nodes near coords style={font=\small\bfseries,color=tinta,anchor=west},
  clip=false]
  \addplot[fill=corB,draw=none,bar shift=0pt] coordinates {(11,1) (4,9) (3,11) (2,17)};
  \addplot[fill=corA,draw=none,bar shift=0pt] coordinates {(7,2) (7,3) (3,10) (2,15) (2,16) (1,19) (1,20)};
  \addplot[fill=corC,draw=none,bar shift=0pt] coordinates {(7,4) (4,7) (4,8)};
  \addplot[fill=cinza!40,draw=none,bar shift=0pt] coordinates {(5,5) (5,6) (3,12) (3,13) (3,14) (2,18)};
\end{axis}
\end{tikzpicture}

\vspace{-1mm}
{\small
\tikz[baseline=-0.6ex]\fill[corA] (0,0) circle (3pt);\ Mecânica\quad
\tikz[baseline=-0.6ex]\fill[corB] (0,0) circle (3pt);\ Eletricidade e magnetismo\quad
\tikz[baseline=-0.6ex]\fill[corC] (0,0) circle (3pt);\ Ondas\quad
\tikz[baseline=-0.6ex]\fill[cinza!40] (0,0) circle (3pt);\ outras áreas}
\end{center}

\subsecao{Ano a ano}
% mapa de calor: fundo mais forte = mais questões naquele ano
\newcommand{\cel}[2]{\cellcolor{#1!\the\numexpr#2*11\relax}#2}
\begin{center}\small
\renewcommand{\arraystretch}{1.35}
\begin{tabular}{l*{5}{>{\centering\arraybackslash}p{1.35cm}}>{\centering\arraybackslash}p{1.5cm}}
\toprule
\textbf{Área} & \textbf{2021} & \textbf{2022} & \textbf{2023} & \textbf{2024} & \textbf{2025} & \textbf{Total}\\
\midrule
Mecânica & \cel{corA}{4} & \cel{corA}{7} & \cel{corA}{4} & \cel{corA}{4} & \cel{corA}{4} & \textbf{23}\\
Eletricidade e magnetismo & \cel{corB}{5} & \cel{corB}{5} & \cel{corB}{2} & \cel{corB}{3} & \cel{corB}{5} & \textbf{20}\\
Ondas & \cel{corC}{3} & \cel{corC}{1} & \cel{corC}{5} & \cel{corC}{3} & \cel{corC}{3} & \textbf{15}\\
Calor (termologia) & \cel{cinza}{4} & \cel{cinza}{1} & \cel{cinza}{4} & \cel{cinza}{4} & \cel{cinza}{2} & \textbf{15}\\
Óptica & \cel{cinza}{0} & \cel{cinza}{1} & \cel{cinza}{0} & \cel{cinza}{1} & \cel{cinza}{1} & \textbf{3}\\
Física moderna & \cel{cinza}{0} & \cel{cinza}{1} & \cel{cinza}{1} & \cel{cinza}{1} & \cel{cinza}{0} & \textbf{3}\\
\midrule
\textbf{Física na prova} & \textbf{16} & \textbf{16} & \textbf{16} & \textbf{16} & \textbf{15} & \textbf{79}\\
\bottomrule
\end{tabular}
\end{center}
\nota{Ondas e calor empatam em 15 questões. Ondas ficou com o 3º lugar porque também aparece como segundo assunto em questões de outras áreas
(18 aparições no total, contra 16 de calor).}

% ---------------------------------------------------------------- EQUAÇÕES E CONCEITOS
\clearpage
\titulo{Equações que resolvem}{As equações que resolvem as questões}
Para cada questão das três áreas campeãs, anotamos quais equações e ideias aparecem na resolução que leva ao gabarito.
A barra mostra em quantas questões cada uma aparece (uma questão pode usar mais de uma).

\begin{tabelaeq}{corA}{\faCarSide\ \ Mecânica · 23 questões}
\linhaeq{\textbf{Desenhar e somar as forças} (peso, normal, atrito, tração)}{todas as questões de leis de Newton começam aqui}{7}
\linhaeq{$v = \dfrac{\Delta s}{\Delta t}$}{viagens, tempos de travessia, parte horizontal dos lançamentos}{6}
\linhaeq{$F_R = m \cdot a$\ \ {\footnotesize($F_R = 0$ com $v$ constante)}}{caminhão freando, polias, balanço}{5}
\linhaeq{$E_c = \dfrac{m v^2}{2}$,\ \ $E_p = m g h$,\ \ $P = \dfrac{E}{\Delta t}$}{potência de motor, painel solar, batida de carro}{5}
\linhaeq{\textbf{Separar o movimento em dois eixos}}{lançamentos, nadador na correnteza, avião com vento}{4}
\linhaeq{\textbf{Atrito:} $F_{\text{at}} = \mu \cdot N$}{caixa na rampa, carga no caminhão, muleta}{4}
\linhaeq{$h = \dfrac{g\,t^2}{2}$,\ \ $v = v_0 - g\,t$}{queda e lançamento vertical ($g = 10$~m/s$^2$)}{3}
\linhaeq{$Q = m \cdot v$,\ \ $F \cdot \Delta t = \Delta Q$}{colisões e proteção em batidas}{2}
\linhaeq{$d = \dfrac{m}{V}$,\ \ $p = \dfrac{F}{A}$}{densidade e pressão: o que flutua, carga sobre a laje}{2}
\end{tabelaeq}

\begin{tabelaeq}{corB}{\faBolt\ \ Eletricidade e magnetismo · 20 questões}
\linhaeq{\textbf{Ler o circuito:} série, paralelo e caminho fechado}{em série a tensão se divide; em paralelo a tensão é a mesma}{10}
\linhaeq{$U = R \cdot i$}{achar corrente, tensão ou resistência em qualquer trecho}{8}
\linhaeq{$P = U \cdot i$,\ \ $P = U^2/R$,\ \ $P = R \cdot i^2$}{aparelhos, disjuntor, brilho de lâmpada, aquecimento}{8}
\linhaeq{\textbf{Pilhas em série somam} (invertida subtrai)}{lanterna, baterias, placas solares}{4}
\linhaeq{$R = \rho\,\dfrac{L}{A}$}{comparar fios: comprido e fino resiste mais}{2}
\linhaeq{$E = P \cdot \Delta t$}{energia em kWh, tempo de recarga, água do chuveiro}{2}
\linhaeq{\textbf{Indução:} campo magnético variando cria corrente}{fogão de indução (2023 e 2025)}{2}
\linhaeq{\textbf{Corrente e campo magnético}}{bússola perto do fio, regra da mão direita}{2}
\linhaeq{$F = q \cdot E$ e linhas de campo elétrico}{tubo de TV antigo, blindagem}{2}
\end{tabelaeq}

\begin{tabelaeq}{corC}{\faWaveSquare\ \ Ondas · 15 questões}
\linhaeq{\textbf{Reconhecer o fenômeno}}{reflexão, difração, interferência, ressonância, Doppler, absorção}{7}
\linhaeq{$v = \lambda \cdot f$\ \ {\footnotesize(luz: $c = \lambda \cdot f$)}}{relacionar frequência e comprimento de onda}{5}
\linhaeq{\textbf{Altura = frequência; intensidade = amplitude}}{grave e agudo, forte e fraco}{4}
\linhaeq{\textbf{Ordem do espectro} (do rádio aos raios gama)}{cores, laser, raios X}{4}
\linhaeq{\textbf{Frequência é da fonte, velocidade é do meio}}{som em barras, na água e no ar}{3}
\linhaeq{$I = \dfrac{P}{A}$ e decibel}{sinal que enfraquece com a distância, ruído}{3}
\linhaeq{$v = \dfrac{\Delta s}{\Delta t}$ para o som {\footnotesize(eco: ida e volta)}}{sonar, prospecção, atraso do som}{2}
\linhaeq{$f = \dfrac{1}{T}$}{ler o período num gráfico}{1}
\end{tabelaeq}

\corcap{tinta}
\nota{As seis equações do kit mínimo ($v = \Delta s/\Delta t$, $F_R = m \cdot a$, $F_{\text{at}} = \mu \cdot N$, $U = R \cdot i$, $P = U \cdot i$ e $v = \lambda \cdot f$)
aparecem na resolução de 33 das 58 questões das três áreas. As outras 25 se resolvem com conceito, leitura de gráfico ou outra equação das tabelas.}

\subsecao{Conta ou conceito? Depende da área}
\begin{center}
\begin{tikzpicture}[x=0.49cm,y=1cm]
  % barras empilhadas: cálculo | conceitual | gráfico ou tabela, com rótulo dentro de cada trecho
  \foreach \y/\nome/\a/\b/\c/\cor in {2/Mecânica/11/9/3/corA, 1/Eletricidade e magnetismo/11/9/0/corB, 0/Ondas/3/7/5/corC}{
    \node[anchor=east,font=\small] at (-0.3,\y) {\nome};
    \fill[\cor] (0,\y-0.3) rectangle (\a,\y+0.3);
    \node[font=\small\bfseries,white] at (\a/2,\y) {\a};
    \fill[\cor!55] (\a+0.08,\y-0.3) rectangle (\a+\b,\y+0.3);
    \node[font=\small\bfseries,text=tinta] at (\a+\b/2,\y) {\b};
    \ifnum\c>0
      \fill[\cor!22] (\a+\b+0.08,\y-0.3) rectangle (\a+\b+\c,\y+0.3);
      \node[font=\small\bfseries,text=tinta] at (\a+\b+\c/2,\y) {\c};
    \fi
  }
  \node[anchor=south,font=\footnotesize\bfseries] at (5.5,2.35) {cálculo};
  \node[anchor=south,font=\footnotesize\bfseries] at (15.5,2.35) {conceito (sem conta)};
  \node[anchor=south west,font=\footnotesize\bfseries] at (20.3,2.35) {gráfico ou tabela};
\end{tikzpicture}
\end{center}
\nota{Em cada barra, o tom escuro é cálculo, o médio é conceito e o claro é leitura de gráfico ou tabela.
Em mecânica e eletricidade, metade das questões pede conta. Em ondas, 12 das 15 são de conceito ou de leitura de gráfico.}

\vspace{3mm}
{\footnotesize\color{cinza}
\textbf{Fontes e método.} Provas e gabaritos oficiais do 2º dia do ENEM, 1ª aplicação, caderno 7 (azul), edições de 2021 a 2025,
publicados pelo INEP em \url{https://download.inep.gov.br/enem/provas_e_gabaritos}. As 45 questões de Ciências da Natureza de cada prova
(números 91 a 135) foram lidas uma a uma. Entraram na contagem as 79 cuja resposta depende principalmente de conteúdo de Física,
classificadas pelo conceito que decide a alternativa correta. Dez delas fazem fronteira com outras disciplinas (oito com Química ou Biologia,
duas com Geografia e Astronomia) e duas foram anuladas pelo INEP (2024, questão 129, e 2025, questão 132). Todas estão na contagem.
Outra classificação razoável pode mudar algum número em uma ou duas unidades, sem alterar os dois primeiros lugares (o 3º está em empate com calor).
As questões reais aparecem resumidas com outras palavras, com ano e número para consulta na prova original. As seis questões de treino são inéditas.\par}
